This work introduced Context Sampling Attention (\cosa{}), a lightweight decoder-side refinement module for remote sensing change detection. The central idea is to use bi-temporal feature correlation as an explicit control signal for refinement, implemented by three coordinated components: correlation-guided sampling, multi-scale context fusion, and learnable residual gating.

On four benchmarks (LEVIR-CD, S2Looking, DSIFN, and CLCD), \cosa{} improves changed-class F1 and IoU over the baseline. Ablations indicate that multi-scale sampling and adaptive gating together are necessary, supporting the view that \cosa{} is a composite mechanism rather than a single attention trick. Cross-backbone tests show strong gains when temporal interaction is limited and near-neutral behavior when fusion is already strong.

Despite these consistent gains, limitations remain. \cosa{} is not uniformly beneficial across architectures. Gains are largest when temporal fusion is relatively weak and can be near-neutral when fusion is already strong. Some test samples degrade, especially in ambiguous or noisy-boundary regions where correlation is unreliable. Learnable gating is validated empirically, but fine-grained analysis of gate dynamics remains limited. Efficiency is discussed via asymptotic cost and measured parameter overhead rather than standardized cross-hardware latency and memory profiling.

These limitations point to several practical future directions. Confidence-aware or uncertainty-aware gating could reduce degradation on difficult samples by suppressing unnecessary refinement. Broader transfer studies across additional transformer and large-scale CD backbones would clarify more precisely where \cosa{} is complementary and where it becomes redundant. Standardized deployment profiling across hardware platforms would also strengthen the practical case for real-world use.
\EOD
